% Options for packages loaded elsewhere
\PassOptionsToPackage{unicode}{hyperref}
\PassOptionsToPackage{hyphens}{url}
\documentclass[
  11pt,
]{article}
\usepackage{xcolor}
\usepackage[margin=25mm]{geometry}
\usepackage{amsmath,amssymb}
\setcounter{secnumdepth}{-\maxdimen} % remove section numbering
\usepackage{iftex}
\ifPDFTeX
  \usepackage[T1]{fontenc}
  \usepackage[utf8]{inputenc}
  \usepackage{textcomp} % provide euro and other symbols
\else % if luatex or xetex
  \usepackage{unicode-math} % this also loads fontspec
  \defaultfontfeatures{Scale=MatchLowercase}
  \defaultfontfeatures[\rmfamily]{Ligatures=TeX,Scale=1}
\fi
\usepackage{lmodern}
\ifPDFTeX\else
  % xetex/luatex font selection
\fi
% Use upquote if available, for straight quotes in verbatim environments
\IfFileExists{upquote.sty}{\usepackage{upquote}}{}
\IfFileExists{microtype.sty}{% use microtype if available
  \usepackage[]{microtype}
  \UseMicrotypeSet[protrusion]{basicmath} % disable protrusion for tt fonts
}{}
\makeatletter
\@ifundefined{KOMAClassName}{% if non-KOMA class
  \IfFileExists{parskip.sty}{%
    \usepackage{parskip}
  }{% else
    \setlength{\parindent}{0pt}
    \setlength{\parskip}{6pt plus 2pt minus 1pt}}
}{% if KOMA class
  \KOMAoptions{parskip=half}}
\makeatother
\setlength{\emergencystretch}{3em} % prevent overfull lines
\providecommand{\tightlist}{%
  \setlength{\itemsep}{0pt}\setlength{\parskip}{0pt}}
\usepackage{bookmark}
\IfFileExists{xurl.sty}{\usepackage{xurl}}{} % add URL line breaks if available
\urlstyle{same}
% fallback for those not using the hyperref driver hyperxmp:
\makeatletter
\@ifundefined{xmpquote}{\newcommand{\xmpquote}[1]{#1}}{}
\makeatother
\hypersetup{
  pdftitle={An exact counterexample to odd-order Hankel H-eigenvalue inheritance},
  pdfauthor={Matthew J. Colbrook},
  hidelinks,
  pdfcreator={LaTeX via pandoc}}

\title{An exact counterexample to odd-order Hankel H-eigenvalue
inheritance}
\author{Matthew J. Colbrook}
\date{11 September 2026}

\usepackage{xurl}
\setlength{\emergencystretch}{3em}
\begin{document}
\maketitle

Department of Applied Mathematics and Theoretical Physics, University of
Cambridge, Cambridge, United Kingdom. Email: m.colbrook@damtp.cam.ac.uk.

This argument was developed with AI assistance at the author's request.
Independent verification is recorded separately; no novelty or external
human peer-review claim is made.

\subsection{Counterexample}\label{counterexample}

Use the canonical TR-15 parameters \(m=3\), \(q=2\) and \(n=2\), and the
common generating vector \[
(h_0,h_1,\ldots,h_6)=(2,0,1,0,2,0,-1).
\] Let \(A\) be the order-three, dimension-three Hankel tensor and \(B\)
the order-six, dimension-two Hankel tensor with entries \[
A_{ijk}=h_{i+j+k-3},\qquad
B_{i_1\cdots i_6}=h_{i_1+\cdots+i_6-6}.
\] Their dimensions satisfy \(\dim A=q(n-1)+1\) and \(\dim B=n\), and
both use all seven entries of the same generating vector.

\textbf{Proposition.} Every real H-eigenvalue of \(A\) is strictly
positive, but \(-1\) is a real H-eigenvalue of \(B\). Consequently the
universal inheritance conjecture in TR-15 is false.

\textbf{Proof.} For a real vector \(x=(x_1,x_2,x_3)\), the first
coordinate of the order-three contraction is \[
(Ax^2)_1=2x_1^2+x_2^2+2x_1x_3+2x_3^2
=(x_1+x_3)^2+x_1^2+x_2^2+x_3^2.
\] This is strictly positive for every nonzero real \(x\). If
\((\lambda,x)\) is any real H-eigenpair, its first equation is
\((Ax^2)_1=\lambda x_1^2\). Thus \(x_1\ne0\) and \(\lambda>0\). In
particular there is no negative H-eigenvalue; no numerical enumeration
of eigenpairs is involved.

For \(v=(0,1)\), the only nonzero product in each component of \(Bv^5\)
has its last five indices equal to two. Hence \[
Bv^5=(h_5,h_6)=(0,-1)=-1\,(v_1^5,v_2^5).
\] This gives the asserted negative H-eigenvalue with an exact nonzero
real eigenvector. All parameters are admissible in the canonical
conjecture, proving its negation. \(\square\)

\subsection{The premise is not
vacuous}\label{the-premise-is-not-vacuous}

Although unnecessary for the counterexample, \(A\) does have real
H-eigenvalues. Its remaining contractions are \[
(Ax^2)_2=2x_1x_2+4x_2x_3,\qquad
(Ax^2)_3=x_1^2+2x_2^2+4x_1x_3-x_3^2.
\] Set \(x=(1,0,t)\) and \(\lambda=2+2t+2t^2\). The second
eigen-equation is automatic and the third is equivalent to \[
2t^4+2t^3+3t^2-4t-1=0.
\] This polynomial is \(-1\) at zero and \(2\) at one, so it has a real
root \(0<t<1\). At that root the displayed \(x,\lambda\) are an
H-eigenpair and \(\lambda>0\).

\subsection{Scope and provenance}\label{scope-and-provenance}

This disproves inheritance from odd lower order in the precise
H-eigenvalue sense. It does not contradict the even-lower-order theorem
or inheritance from a positive-semidefinite associated Hankel matrix.
Indeed the associated \(4\times4\) Hankel matrix has bottom-right entry
\(h_6=-1\), so that stronger positive-semidefinite hypothesis fails. The
tensor \(B\) also takes the value \(-1\) on the polynomial evaluation at
\((0,1)\), consistently with its negative H-eigenvalue.

The prior incomplete tensor archives contained only a qualitative
counterexample summary, with no generating vector or proof. The explicit
coefficients and proof above were derived anew in this research pass; no
absent manuscript or code is represented as recovered.

\subsection{References}\label{references}

\begin{enumerate}
\def\labelenumi{\arabic{enumi}.}
\tightlist
\item
  W. Ding, L. Qi and Y. Wei, \emph{Inheritance properties and
  sum-of-squares decomposition of Hankel tensors: theory and
  algorithms}, BIT Numerical Mathematics 57 (2017), 169-190, final
  Section 4, concluding odd-order inheritance conjecture; Section 2
  fixes the common-vector conventions.
  \href{https://www.polyu.edu.hk/ama/staff/new/qilq/BIT-DQW.pdf}{Author's
  journal PDF}, \href{https://doi.org/10.1007/s10543-016-0622-0}{DOI}.
\item
  \href{https://github.com/ajt60gaibb/OpenProblemsInNLA/blob/main/tensor-computations/TR-15/README.md}{Canonical
  TR-15 statement}, checked 11 September 2026.
\end{enumerate}

\end{document}
